\documentclass[../DoAn.tex]{subfiles}
\begin{document}

Phụ lục này bổ sung một số use case chưa được đặc tả chi tiết trong Chương \ref{chapter:Requirements} nhưng vẫn quan trọng đối với vận hành hệ thống.

\section{Đặc tả use case Quản lý wishlist}
\label{appendix:B.1}

{\normalsize
\setlength{\LTleft}{\fill}
\setlength{\LTright}{\fill}
\setlength{\tabcolsep}{5pt}
\renewcommand{\arraystretch}{1.25}
\begin{longtable}{|p{3.2cm}|p{10.4cm}|}
\caption{Đặc tả use case Quản lý wishlist}\\
\hline
\textbf{Mục} & \textbf{Nội dung} \\ \hline
Tên use case & Quản lý wishlist \\ \hline
Tác nhân chính & Khách hàng \\ \hline
Mục tiêu & Cho phép khách hàng lưu sản phẩm quan tâm để quay lại mua hoặc so sánh sau. \\ \hline
Tiền điều kiện & Khách hàng đã đăng nhập; sản phẩm còn tồn tại trong catalog. \\ \hline
Luồng sự kiện chính &
\begin{minipage}[t]{10.1cm}
\begin{enumerate}[leftmargin=*,topsep=0pt,itemsep=0pt]
    \item Khách hàng bấm lưu sản phẩm từ trang shop hoặc chi tiết sản phẩm.
    \item Frontend gửi product id tới API wishlist.
    \item Backend kiểm tra sản phẩm tồn tại, người dùng hiện tại và ràng buộc unique theo cặp user/product.
    \item Hệ thống lưu hoặc trả về trạng thái đã có trong wishlist.
    \item Khách hàng mở trang wishlist để xem và xóa sản phẩm đã lưu.
\end{enumerate}
\end{minipage} \\ \hline
Hậu điều kiện & Bảng \texttt{wishlist\_items} phản ánh đúng danh sách sản phẩm người dùng đã lưu. \\ \hline
\end{longtable}
}

\section{Đặc tả use case Đánh giá sản phẩm sau mua}
\label{appendix:B.2}

{\normalsize
\setlength{\LTleft}{\fill}
\setlength{\LTright}{\fill}
\setlength{\tabcolsep}{5pt}
\renewcommand{\arraystretch}{1.25}
\begin{longtable}{|p{3.2cm}|p{10.4cm}|}
\caption{Đặc tả use case Đánh giá sản phẩm sau mua}\\
\hline
\textbf{Mục} & \textbf{Nội dung} \\ \hline
Tên use case & Đánh giá sản phẩm sau mua \\ \hline
Tác nhân chính & Khách hàng \\ \hline
Mục tiêu & Thu thập đánh giá thực tế từ khách đã mua để tăng độ tin cậy của catalog và hỗ trợ gợi ý sản phẩm. \\ \hline
Tiền điều kiện & Khách hàng đã đăng nhập và có đơn hàng hợp lệ chứa sản phẩm cần đánh giá. \\ \hline
Luồng sự kiện chính &
\begin{minipage}[t]{10.1cm}
\begin{enumerate}[leftmargin=*,topsep=0pt,itemsep=0pt]
    \item Khách hàng mở phần đánh giá ở trang chi tiết sản phẩm.
    \item Frontend gọi API kiểm tra người dùng có quyền đánh giá hay không.
    \item Khách hàng nhập số sao và bình luận.
    \item Backend validate số sao từ 1 đến 5 và kiểm tra ràng buộc mỗi user chỉ review một sản phẩm một lần.
    \item Hệ thống lưu review và cập nhật thống kê rating của sản phẩm.
\end{enumerate}
\end{minipage} \\ \hline
Hậu điều kiện & Review được lưu trong bảng \texttt{reviews}; sản phẩm hiển thị điểm trung bình và số lượng đánh giá mới. \\ \hline
\end{longtable}
}

\section{Đặc tả use case Yêu cầu đổi trả}
\label{appendix:B.3}

{\normalsize
\setlength{\LTleft}{\fill}
\setlength{\LTright}{\fill}
\setlength{\tabcolsep}{5pt}
\renewcommand{\arraystretch}{1.25}
\begin{longtable}{|p{3.2cm}|p{10.4cm}|}
\caption{Đặc tả use case Yêu cầu đổi trả}\\
\hline
\textbf{Mục} & \textbf{Nội dung} \\ \hline
Tên use case & Yêu cầu đổi trả \\ \hline
Tác nhân chính & Khách hàng, nhân viên quản trị đơn hàng \\ \hline
Mục tiêu & Cho phép khách hàng gửi yêu cầu đổi trả và nhân viên xử lý theo trạng thái nghiệp vụ. \\ \hline
Tiền điều kiện & Đơn hàng thuộc về khách hàng, có sản phẩm đủ điều kiện đổi trả theo chính sách. \\ \hline
Luồng sự kiện chính &
\begin{minipage}[t]{10.1cm}
\begin{enumerate}[leftmargin=*,topsep=0pt,itemsep=0pt]
    \item Khách hàng mở chi tiết đơn hàng và chọn yêu cầu đổi trả.
    \item Khách hàng nhập lý do và số lượng từng dòng hàng cần đổi trả.
    \item Backend kiểm tra quyền sở hữu đơn, số lượng hợp lệ và trạng thái đơn.
    \item Hệ thống tạo bản ghi \texttt{order\_returns} và \texttt{order\_return\_items}.
    \item Nhân viên admin xem danh sách yêu cầu, duyệt, hoàn tất hoặc từ chối kèm ghi chú.
\end{enumerate}
\end{minipage} \\ \hline
Hậu điều kiện & Yêu cầu đổi trả có trạng thái rõ ràng để khách hàng và nhân viên theo dõi. \\ \hline
\end{longtable}
}

\section{Đặc tả use case Quản lý embedding}
\label{appendix:B.4}

{\normalsize
\setlength{\LTleft}{\fill}
\setlength{\LTright}{\fill}
\setlength{\tabcolsep}{5pt}
\renewcommand{\arraystretch}{1.25}
\begin{longtable}{|p{3.2cm}|p{10.4cm}|}
\caption{Đặc tả use case Quản lý embedding}\\
\hline
\textbf{Mục} & \textbf{Nội dung} \\ \hline
Tên use case & Quản lý embedding \\ \hline
Tác nhân chính & Quản trị viên, nhân viên nội dung \\ \hline
Mục tiêu & Cho phép nhân sự vận hành reindex sản phẩm hoặc kho tri thức khi dữ liệu thay đổi. \\ \hline
Tiền điều kiện & Người dùng có quyền phù hợp; backend được cấu hình OpenAI API key; PostgreSQL/pgvector hoạt động. \\ \hline
Luồng sự kiện chính &
\begin{minipage}[t]{10.1cm}
\begin{enumerate}[leftmargin=*,topsep=0pt,itemsep=0pt]
    \item Nhân viên mở trang quản lý embeddings trong admin.
    \item Nhân viên chọn collection sản phẩm hoặc knowledge.
    \item Backend đọc dữ liệu đang active, tạo document và gọi PGVector store để thêm embedding.
    \item Với knowledge, hệ thống đánh giá lại forum thread đủ điều kiện trước khi index.
    \item Hệ thống trả về số lượng document/chunk đã reindex.
\end{enumerate}
\end{minipage} \\ \hline
Hậu điều kiện & Collection \texttt{petshop\_products} hoặc \texttt{petshop\_knowledge} được đồng bộ lại với dữ liệu nghiệp vụ. \\ \hline
\end{longtable}
}

\end{document}
